\chapter{Introduction and Counting Basics}

Math 239 has two halves. The first half is \emph{enumeration}: given a set of objects
described by some rule, how many objects of each ``size'' are there? The second half is an
introduction to \emph{graph theory}. This chapter collects the basic counting tools that the
rest of the enumeration half is built on: sums and products of sets, binomial coefficients,
bijections and combinatorial proofs.

\section{Sets, sums and products}

All sets in this course are finite unless stated otherwise. For a finite set $S$ we write
$\abs{S}$ for its number of elements, and $[n] = \{1, 2, \dots, n\}$.

\begin{definition}[Disjoint union, Cartesian product]
Sets $A$ and $B$ are \emph{disjoint} if $A \cap B = \emptyset$; we then call $A \cup B$ a
\emph{disjoint union}. The \emph{Cartesian product} of $A$ and $B$ is
$A \times B = \{(a,b) : a \in A,\ b \in B\}$, and more generally
$A_1 \times \cdots \times A_k$ is the set of $k$-tuples $(a_1,\dots,a_k)$ with $a_i \in A_i$.
We write $A^k$ for $A \times \cdots \times A$ ($k$ factors).
\end{definition}

\begin{theorem}[Sum and product rules]
\begin{enumerate}
  \item If $A$ and $B$ are disjoint, then $\abs{A \cup B} = \abs{A} + \abs{B}$. In general,
        $\abs{A\cup B} = \abs{A} + \abs{B} - \abs{A \cap B}$.
  \item $\abs{A \times B} = \abs{A}\cdot\abs{B}$, and so $\abs{A^k} = \abs{A}^k$.
\end{enumerate}
\end{theorem}

\begin{example}
The number of subsets of $[n]$ is $2^n$: a subset $S$ is determined by the $n$-tuple
$(s_1,\dots,s_n) \in \{0,1\}^n$ where $s_i = 1$ iff $i \in S$, and $\abs{\{0,1\}^n} = 2^n$.
\end{example}

\section{Permutations and binomial coefficients}

\begin{definition}[Permutation]
A \emph{permutation} of a set $S$ with $n$ elements is an ordered list of the elements of $S$
in which every element appears exactly once. More generally a $k$-permutation is an ordered
list of $k$ distinct elements of $S$.
\end{definition}

\begin{proposition}
The number of permutations of an $n$-set is $n! = n(n-1)\cdots 1$, and the number of
$k$-permutations is $n(n-1)\cdots(n-k+1) = \dfrac{n!}{(n-k)!}$ for $0 \le k \le n$.
\end{proposition}
\begin{proof}
There are $n$ choices for the first entry, then $n-1$ for the second (anything not yet used),
and so on; after $k$ entries we have $n(n-1)\cdots(n-k+1)$ lists.
\end{proof}

\begin{definition}[Binomial coefficient]
For integers $n, k \ge 0$, $\binom{n}{k}$ is the number of $k$-element subsets of an $n$-element
set. (So $\binom{n}{k} = 0$ when $k > n$.)
\end{definition}

\begin{theorem}
For $0 \le k \le n$,
\[
\binom{n}{k} = \frac{n!}{k!\,(n-k)!} = \frac{n(n-1)\cdots(n-k+1)}{k!}.
\]
\end{theorem}
\begin{proof}
Count $k$-permutations of $[n]$ in two ways. Directly there are $n!/(n-k)!$. Alternatively,
first choose the set of $k$ elements that appear ($\binom{n}{k}$ ways), then order them
($k!$ ways). Hence $\binom{n}{k}\, k! = n!/(n-k)!$.
\end{proof}

The second formula $\frac{n(n-1)\cdots(n-k+1)}{k!}$ makes sense for \emph{any} $n$ (even
negative or non-integer) and we use it to \emph{define} $\binom{n}{k}$ in general; this is
needed later for the negative binomial theorem.

\begin{theorem}[Binomial theorem]
For $n \ge 0$, $\displaystyle (1+x)^n = \sum_{k=0}^{n} \binom{n}{k} x^k$.
\end{theorem}
\begin{proof}
Expand $(1+x)(1+x)\cdots(1+x)$: each term of the expansion picks either $1$ or $x$ from each of
the $n$ factors. The term $x^k$ arises once for each choice of the $k$ factors that contribute
$x$, i.e.\ $\binom{n}{k}$ times.
\end{proof}

\section{Bijections}

\begin{definition}[Bijection]
A function $f : A \to B$ is a \emph{bijection} (one-to-one correspondence) if it is both
injective ($f(a) = f(a') \Rightarrow a = a'$) and surjective (every $b \in B$ is $f(a)$ for some
$a$). Equivalently, $f$ has an inverse $g : B \to A$ with $g\circ f = \mathrm{id}_A$ and
$f \circ g = \mathrm{id}_B$.
\end{definition}

\begin{theorem}
For finite sets, $\abs{A} = \abs{B}$ if and only if there exists a bijection $A \to B$.
\end{theorem}

In practice, to prove a map is a bijection it is usually easiest to write down the candidate
inverse and check both compositions are identities.

\begin{example}
Complementation $S \mapsto [n] \setminus S$ is a bijection from $k$-subsets of $[n]$ to
$(n-k)$-subsets of $[n]$ (it is its own inverse). Hence $\binom{n}{k} = \binom{n}{n-k}$.
\end{example}

\begin{example}
Let $A$ be the set of subsets of $[n]$ of even size and $B$ the set of subsets of odd size,
$n \ge 1$. The map $f(S) = S \triangle \{1\}$ (add $1$ if absent, remove it if present) changes
the size by exactly one, so it maps $A$ to $B$, and $f\circ f = \mathrm{id}$. So $\abs{A} = \abs{B}
= 2^{n-1}$, which gives $\sum_k (-1)^k \binom{n}{k} = 0$.
\end{example}

\section{Combinatorial proofs}

A \emph{combinatorial proof} of an identity $L = R$ finds a set $S$ and counts it in two
different ways, one giving $L$ and the other $R$ (or exhibits a bijection between a set of
size $L$ and a set of size $R$). Such proofs usually explain \emph{why} an identity holds better
than algebra does.

\begin{theorem}[Pascal's identity]
For $n, k \ge 1$, $\displaystyle \binom{n}{k} = \binom{n-1}{k-1} + \binom{n-1}{k}$.
\end{theorem}
\begin{proof}
Count $k$-subsets of $[n]$. Those containing $n$ correspond to $(k-1)$-subsets of $[n-1]$
(remove $n$); those not containing $n$ are exactly the $k$-subsets of $[n-1]$. The two kinds
are disjoint, so the sum rule gives the identity.
\end{proof}

\begin{theorem}
$\displaystyle \sum_{k=0}^n \binom{n}{k} = 2^n$.
\end{theorem}
\begin{proof}
Both sides count all subsets of $[n]$: the left side sorts them by size.
\end{proof}

\begin{theorem}
For $n \ge 1$, $\displaystyle \sum_{k=0}^n k\binom{n}{k} = n2^{n-1}$.
\end{theorem}
\begin{proof}
Count pairs $(S, x)$ where $S \subseteq [n]$ and $x \in S$ (a ``committee with a chair'').
Choosing $S$ first with $\abs{S} = k$ and then $x$ gives $\sum_k k\binom{n}{k}$. Choosing $x$
first ($n$ ways) and then any subset of the remaining $n-1$ elements to join it ($2^{n-1}$ ways)
gives $n2^{n-1}$.
\end{proof}

\begin{theorem}[Vandermonde's identity]
$\displaystyle \binom{m+n}{k} = \sum_{i=0}^{k} \binom{m}{i}\binom{n}{k-i}$.
\end{theorem}
\begin{proof}
Choose $k$ people from a group of $m$ women and $n$ men. Sorting the choices by the number $i$
of women chosen, there are $\binom{m}{i}\binom{n}{k-i}$ choices with exactly $i$ women.
\end{proof}

\begin{theorem}[Hockey stick]
$\displaystyle \sum_{i=k}^{n} \binom{i}{k} = \binom{n+1}{k+1}$.
\end{theorem}
\begin{proof}
Count $(k+1)$-subsets of $[n+1]$ by their largest element. If the largest element is $i+1$,
the other $k$ elements form a $k$-subset of $[i]$, which can be chosen in $\binom{i}{k}$ ways,
and $i$ ranges from $k$ to $n$.
\end{proof}

\begin{important}[Writing a combinatorial proof]
State clearly (1) the set being counted, (2) the first way of counting it, and (3) the second
way. Each way must count every element exactly once. An argument of the form ``the left side
counts this, and the right side obviously counts the same thing'' without justification of the
second count is not a proof.
\end{important}

\section{Multisets}

\begin{definition}[Multiset]
A \emph{multiset} of size $n$ with elements of $t$ types is a selection of $n$ objects where
objects of the same type are indistinguishable and repetition is allowed. Equivalently it is a
$t$-tuple $(m_1, \dots, m_t)$ of non-negative integers with $m_1 + \cdots + m_t = n$
($m_i$ = number of objects of type $i$).
\end{definition}

\begin{theorem}
The number of multisets of size $n$ with $t \ge 1$ types is
$\displaystyle \binom{n+t-1}{t-1}$.
\end{theorem}
\begin{proof}[Proof (stars and bars)]
Represent $(m_1,\dots,m_t)$ as a row of $n$ stars and $t-1$ bars: $m_1$ stars, a bar, $m_2$
stars, a bar, \dots, $m_t$ stars. This is a bijection between multisets and arrangements of
$n+t-1$ symbols of which $t-1$ are bars, and such an arrangement is determined by the positions
of the bars: $\binom{n+t-1}{t-1}$ choices.
\end{proof}

\begin{example}
The number of non-negative integer solutions of $a+b+c = 7$ is $\binom{7+2}{2} = 36$. The
number of \emph{positive} solutions is found by substituting $a' = a-1$ etc.:
$a'+b'+c' = 4$ has $\binom{6}{2} = 15$ non-negative solutions.
\end{example}

\section{Summary}
\begin{itemize}
  \item Sum rule for disjoint unions, product rule for Cartesian products; $2^n$ subsets of $[n]$.
  \item $n!$ permutations, $\binom{n}{k} = \frac{n!}{k!(n-k)!}$ subsets of size $k$; binomial theorem.
  \item Two finite sets have the same size iff there is a bijection between them.
  \item Combinatorial proofs: count one set two ways (Pascal, Vandermonde, hockey stick,
        $\sum k\binom{n}{k} = n2^{n-1}$).
  \item Multisets of size $n$ from $t$ types: $\binom{n+t-1}{t-1}$ (stars and bars).
\end{itemize}
